% sec:sr-causal — Light Cones, Causal Structure, and Null Vectors
\section{Light Cones, Causal Structure, and Null Vectors}
\label{sec:sr-causal}

A ray tracer needs to answer one question for every pixel: given a camera at
some event in spacetime, which direction does this pixel correspond to, and
where does a photon from that direction originate?  The answer lives in the set
of null vectors at the camera's location---the directions along which light can
travel.  This set has a rich geometric structure: it forms a cone in the tangent
space, and the spatial directions on that cone are in one-to-one correspondence
with points on the observer's sky.  This section makes that correspondence
precise.

\paragraph{The light cone.}

At any event $p$ in Minkowski spacetime, the \emph{light cone} is the set of
all null directions---nonzero vectors $k^\mu$ satisfying
%
\begin{equation}
  \eta_{\mu\nu}\, k^\mu\, k^\nu = 0 \,.
  \label{eq:null-condition}
\end{equation}
%
In \cref{sec:sr-minkowski} we classified individual vectors as timelike, null,
or spacelike.  The light cone is the \emph{boundary} that separates the
timelike region from the spacelike region: away from the vertex $k^\mu = 0$,
it is a three-dimensional hypersurface in the four-dimensional tangent space
at $p$.

Expanding \cref{eq:null-condition} in Cartesian coordinates gives
%
\begin{equation}
  -(k^0)^2 + (k^1)^2 + (k^2)^2 + (k^3)^2 = 0 \,,
  \label{eq:null-expanded}
\end{equation}
%
which rearranges to $(k^0)^2 = \abs{\mathbf{k}}^2$, where $\abs{\mathbf{k}}^2
= (k^1)^2 + (k^2)^2 + (k^3)^2$ is the squared Euclidean norm of the spatial
part.  This has two solutions: $k^0 = +\abs{\mathbf{k}}$ (future-directed) and
$k^0 = -\abs{\mathbf{k}}$ (past-directed).  The light cone therefore splits
into two halves:
%
\begin{itemize}
  \item The \textbf{future light cone}: all null vectors with $k^0 > 0$.
    These point toward the future---photons emitted from $p$ travel along
    these directions.
  \item The \textbf{past light cone}: all null vectors with $k^0 < 0$.
    These point toward the past---photons arriving at $p$ came from
    these directions.
\end{itemize}
%
The ray tracer works with the past light cone: it traces photons backward in
time from the camera to their sources, because that is far more efficient than
tracing photons forward from every source and hoping some reach the camera.
\physnote{Backward ray tracing is standard in computer graphics
\cite{Marschner2016}; the relativistic version merely replaces straight lines
with null geodesics.}

\begin{figure}[htbp]
  \centering
  % light-cone-causal.tex — 2+1 light cone with causal regions
% Inputable TikZ fragment for fig:light-cone-causal
% Requires: tikz with calc, arrows.meta, decorations.markings (loaded by ehbook.cls)
%
\begin{tikzpicture}[
    >={Stealth[length=5pt]},
    font=\footnotesize,
  ]

  \def\h{2.2}     % cone height
  \def\rx{1.65}   % ellipse semi-major (x direction)
  \def\ry{0.55}   % ellipse semi-minor (perspective depth)

  % Tangent points: lines from apex (0,0) tangent to ellipse at height h.
  %   tx = rx·sqrt(h²-ry²)/h,  ty = (h²-ry²)/h,  tangle = acos(tx/rx)
  \pgfmathsetmacro{\tx}{\rx*sqrt(\h*\h - \ry*\ry)/\h}
  \pgfmathsetmacro{\ty}{(\h*\h - \ry*\ry)/\h}
  \pgfmathsetmacro{\tangle}{acos(\tx/\rx)}

  % ── Past cone ─────────────────────────────────────────────────────

  % Back arc of past ellipse (dashed)
  \draw[EHMarginGray, thin, densely dashed]
    (-\tx,-\ty) arc ({180-\tangle}:{360+\tangle}:{\rx} and {\ry});

  % Past cone fill (very light)
  \fill[EHJetBlue, fill opacity=0.05]
    (0,0) -- (\tx,-\ty) arc ({\tangle}:{180-\tangle}:{\rx} and {\ry}) -- cycle;

  % Past cone surface lines
  \draw[EHMarginGray, thin] (0,0) -- (-\tx,-\ty);
  \draw[EHMarginGray, thin] (0,0) -- ( \tx,-\ty);

  % Front arc of past ellipse
  \draw[EHMarginGray, thin]
    (\tx,-\ty) arc ({\tangle}:{180-\tangle}:{\rx} and {\ry});

  % ── Axes (behind future cone) ─────────────────────────────────────

  \draw[->, thick, EHMarginGray] (-2.6,0) -- (2.6,0)
    node[right, font=\small] {$x$};
  \draw[->, thick, EHMarginGray] (0,0) -- (-0.85,-0.5)
    node[below left, font=\small] {$y$};

  % ── Future cone ───────────────────────────────────────────────────

  % Back arc of future ellipse (dashed)
  \draw[EHHorizonCrimson!50, thin, densely dashed]
    (\tx,\ty) arc ({-\tangle}:{180+\tangle}:{\rx} and {\ry});

  % Future cone interior fill — I^+(p)
  \fill[EHAccretionGold, fill opacity=0.10]
    (0,0) -- (-\tx,\ty)
      arc ({180+\tangle}:{360-\tangle}:{\rx} and {\ry}) -- cycle;

  % Cone surface lines — boundary of J^+(p)
  \draw[EHHorizonCrimson, thick] (0,0) -- (-\tx,\ty);
  \draw[EHHorizonCrimson, thick] (0,0) -- ( \tx,\ty);

  % Front arc of future ellipse
  \draw[EHHorizonCrimson, thick]
    (-\tx,\ty) arc ({180+\tangle}:{360-\tangle}:{\rx} and {\ry});

  % ── t-axis (drawn last so it's on top) ────────────────────────────

  \draw[->, thick, EHMarginGray] (0,-2.8) -- (0,2.8)
    node[above, font=\small] {$t$};

  % ── Event p ────────────────────────────────────────────────────────

  \fill (0,0) circle (1.5pt);
  \node[below right=1pt, font=\small] at (0,0) {$p$};

  % ── Timelike worldline (gold, through cone interior) ──────────────

  \draw[very thick, EHAccretionGold,
        decoration={markings, mark=at position 0.72 with {\arrow{>}}},
        postaction={decorate}]
    (0.15,-1.5) .. controls (0.05,-0.5) and (-0.1,0.5) .. (-0.25,1.8);

  % ── Spacelike separation (blue, outside cone) ─────────────────────

  \draw[->, very thick, EHJetBlue]
    (0,0) -- (2.1,0.2);

  % ── Labels ─────────────────────────────────────────────────────────

  \node[text=EHAccretionGold!85!black] at (0.4,1.15) {$I^+(p)$};
  \node[text=EHHorizonCrimson!85!black, right=1pt] at (\tx,\ty) {$J^+(p)$};
  \node[text=EHJetBlue!85!black, above=1pt] at (2.1,0.2) {elsewhere};

\end{tikzpicture}

  \caption[The light cone at an event $p$ in Minkowski spacetime]{The light cone at an event $p$ in Minkowski spacetime.  The
    chronological future $I^+(p)$ (gold interior) contains all events
    reachable by timelike curves---massive particles.  The causal future
    $J^+(p)$ adds the cone boundary itself (crimson), reached by null
    geodesics---photons.  Events outside both cones are spacelike-separated
    from $p$: no signal can connect them.  The ray tracer works with the past
    light cone, tracing photons backward from the camera to their sources.}
  \label{fig:light-cone-causal}
\end{figure}

\paragraph{Causal structure.}

The light cone at an event $p$ divides the tangent space---and, by extension,
the rest of spacetime---into causally distinct regions.  The \emph{chronological
future} $I^+(p)$ is the set of events reachable from $p$ by timelike curves
(massive particles).  The \emph{causal future} $J^+(p)$ includes both timelike
and null curves---it adds the events connected to $p$ by light signals.  The
boundary of $J^+(p)$ is precisely the future light cone of $p$.  Analogous
definitions hold for the past: $I^-(p)$ and $J^-(p)$.

Events outside $J^+(p) \cup J^-(p)$ are \emph{causally disconnected} from $p$:
no signal, massive or massless, can travel between them.  The separation
between $p$ and any such event is spacelike.  This is the geometric content of
the statement that nothing travels faster than light---it is not a speed limit
imposed from outside, but a consequence of the cone structure built into the
metric.
\histnote{Penrose's conformal diagrams \cite{Penrose1963} make the global
causal structure of a spacetime visible at a glance.  We will not need them
until \cref{ch:schwarzschild}, where the causal structure of the black hole
interior becomes central.}

\paragraph{Null vectors and the observer's sky.}

The null condition \cref{eq:null-condition} reduces a vector from four
components down to three free parameters.  But there is a further redundancy: if $k^\mu$ is
null, so is $\lambda k^\mu$ for any $\lambda > 0$.  Rescaling a photon's
four-momentum changes its energy but not its direction.  The physically
distinct null directions therefore form a \emph{two}-parameter family---they
are labelled by a point on a two-sphere.

To see this explicitly, write a future-directed null vector as
%
\begin{equation}
  k^\mu = E\,(1,\, \sin\theta\cos\varphi,\,
               \sin\theta\sin\varphi,\, \cos\theta) \,,
  \label{eq:null-parametrisation}
\end{equation}
%
where $E > 0$ is the photon energy (the freedom we quotient out) and $(\theta,
\varphi)$ are the standard spherical angles on $S^2$.  The two angles label the
directions on the observer's \emph{celestial sphere}---the sky.  Each pixel in
a rendered image corresponds to a specific $(\theta, \varphi)$, and therefore
to a specific null direction $k^\mu$.  Substituting into
\cref{eq:null-condition} confirms the null condition: $-E^2 + E^2(\sin^2\theta
\cos^2\varphi + \sin^2\theta\sin^2\varphi + \cos^2\theta) = -E^2 + E^2 = 0$.

This is the starting point for every ray trace: pick a pixel, compute
the null vector, and propagate it through spacetime.
\xrefnote{The camera model in \cref{ch:camera-rendering} maps pixel
coordinates $(i, j)$ to angles $(\theta, \varphi)$ using the observer's
tetrad from \cref{sec:sr-lorentz}.}

\paragraph{The null covector and momentum.}

Lowering the index on a null vector yields a covector, the photon's
four-momentum $p_\mu$:
%
\begin{equation}
  p_\mu = \eta_{\mu\nu}\, k^\nu \,.
  \label{eq:photon-momentum}
\end{equation}
%
With signature $(-,+,+,+)$, the time component changes sign:
$p_0 = -k^0 = -E$, while the spatial components are unchanged in Cartesian
coordinates: $p_i = k^i$.  The null condition expressed in terms of the
covector is $\eta^{\mu\nu}\,p_\mu\,p_\nu = 0$, which reads
%
\begin{equation}
  -(p_0)^2 + (p_1)^2 + (p_2)^2 + (p_3)^2 = 0 \,,
  \label{eq:null-covector}
\end{equation}
%
or equivalently $p_0 = -\abs{\mathbf{p}}$, where the sign reflects our choice
that $p_0 < 0$ for a future-directed photon.  The covector form is the natural
variable for the ray tracer: the time component $p_0 = -E$ is minus the photon
energy as measured by a static observer, and the spatial components encode its
direction of travel.
\physnote{The minus sign on $p_0$ is a perennial source of sign errors.  In our
$(-,+,+,+)$ convention, a physical photon has $p_0 < 0$ and $E = -p_0 > 0$.}

The Hamiltonian ray tracer integrates the geodesic equations in terms of
$p_\mu$ because the covariant components have simpler evolution equations in
the presence of symmetry---along a geodesic, a Killing vector $\xi^\mu$ gives a
conserved quantity $p_\mu\,\xi^\mu$ directly, without further metric
contraction.

\paragraph{A worked example: constructing a photon's initial data.}

Suppose a camera at rest in Minkowski spacetime (so its tetrad is the identity)
needs to trace a ray toward the direction $(\theta, \varphi) = (60°, 45°)$.
The relevant trigonometric values are $\sin 60° = \sqrt{3}/2$,
$\cos 60° = 1/2$, $\cos 45° = \sin 45° = 1/\sqrt{2}$, and
$\sin 60°\cos 45° = \sqrt{3}/(2\sqrt{2}) = \sqrt{6}/4$.  The initial null
vector is therefore
%
\begin{equation}
  k^\mu = E\bigl(1,\, \tfrac{\sqrt{6}}{4},\,
                     \tfrac{\sqrt{6}}{4},\, \tfrac{1}{2}\bigr) \,.
  \label{eq:photon-example}
\end{equation}
%
The covector is $p_\mu = (-E,\, E\sqrt{6}/4,\,
E\sqrt{6}/4,\, E/2)$.  We can verify the null condition:
$\eta^{\mu\nu} p_\mu p_\nu = -E^2 + 6E^2/16 + 6E^2/16 + E^2/4
= -E^2 + 3E^2/8 + 3E^2/8 + 2E^2/8 = -E^2 + E^2 = 0$.

The energy $E$ drops out of the geodesic equations for null rays (the
trajectory depends only on direction, not on energy), so we can set $E = 1$
without loss of generality.  This gives the concrete initial four-momentum
$p_\mu = (-1,\, \sqrt{6}/4,\, \sqrt{6}/4,\, 1/2)$ that the integrator takes as
input.
\warnnote{Setting $E = 1$ is a choice of affine parametrisation, not a physical
statement about the photon's energy.  Observables like redshift involve
\emph{ratios} of energies measured by different observers, so the absolute
normalisation is irrelevant.}

\paragraph{From flat to curved: the sky in any spacetime.}

In Minkowski spacetime, the null parametrisation \cref{eq:null-parametrisation}
uses the coordinate axes as the reference frame.  For a boosted observer, we
replace the coordinate basis with the observer's tetrad from
\cref{sec:sr-lorentz}: the spatial directions
$(\sin\theta\cos\varphi,\, \sin\theta\sin\varphi,\, \cos\theta)$
are expanded in the tetrad vectors $e^\mu{}_{(1)}, e^\mu{}_{(2)},
e^\mu{}_{(3)}$ instead of the coordinate vectors, and the time direction uses
$e^\mu{}_{(0)}$.  The null vector for a pixel at $(\theta, \varphi)$ as
measured by observer $\mathcal{O}$ becomes
%
\begin{equation}
  k^\mu = E\Bigl(e^\mu{}_{(0)}
    + \sin\theta\cos\varphi\; e^\mu{}_{(1)}
    + \sin\theta\sin\varphi\; e^\mu{}_{(2)}
    + \cos\theta\; e^\mu{}_{(3)}\Bigr) \,.
  \label{eq:null-tetrad}
\end{equation}
%
This expression is valid in \emph{any} spacetime: flat or curved, static or
dynamic.  The tetrad encodes the observer's state of motion; the angles encode
the pixel direction; the metric enters only through the tetrad orthonormality
condition.  To confirm that \cref{eq:null-tetrad} is indeed null, substitute
into $\eta_{\mu\nu}\,k^\mu\,k^\nu$ and use the orthonormality condition
$\eta_{\mu\nu}\,e^\mu{}_{(a)}\,e^\nu{}_{(b)} = \eta_{ab}$: the cross terms
vanish, the timelike leg contributes $-E^2$, and the three spatial legs
contribute $E^2(\sin^2\theta\cos^2\varphi + \sin^2\theta\sin^2\varphi +
\cos^2\theta) = E^2$, giving $-E^2 + E^2 = 0$.
\implnote{The camera module constructs $k^\mu$ from a pixel grid and the
observer's tetrad, then hands the resulting null vector to the geodesic
integrator.}

The chapter introduction promised that this parametrisation of null
vectors would reappear verbatim in the camera model, and
\cref{eq:null-tetrad} is that parametrisation.

In curved spacetime, the tetrad vectors $e^\mu{}_{(a)}$ become functions of
position---they satisfy
$\metric\,e^\mu{}_{(a)}\,e^\nu{}_{(b)} = \eta_{ab}$ rather than the Minkowski
version---but \cref{eq:null-tetrad} itself is unchanged.  The flat-spacetime
machinery of this chapter is the complete template for relativistic ray tracing;
what remains is to let the metric vary.
\xrefnote{The geodesic equation, derived in
\cref{ch:differential-geometry}, governs how $k^\mu$ evolves as the ray
propagates through curved spacetime.}
